\mysection{strstr()の例}
\label{strstr_example}
\myindex{\CStandardLibrary!strstr()}

GCCが時々文字列の一部を使用できるという事実に戻りましょう：\myref{use_parts_of_C_strings}。

\IT{strstr()} \CCpp標準ライブラリ関数は、文字列内の任意の出現を見つけるために使用されます。
これが私たちが行うことです：

\begin{lstlisting}[style=customc]
#include <string.h>
#include <stdio.h>

int main()
{
	char *s="Hello, world!";
	char *w=strstr(s, "world");

	printf ("%p, [%s]\n", s, s);
	printf ("%p, [%s]\n", w, w);
};
\end{lstlisting}

出力は：

\begin{lstlisting}
0x8048530, [Hello, world!]
0x8048537, [world!]
\end{lstlisting}

元の文字列のアドレスと\IT{strstr()}が返した部分文字列のアドレスの差は7です。
実際、\q{Hello, }文字列の長さは7文字です。

2回目の呼び出し時に\printf{}関数は、渡された文字列の前に他の文字があることを知らず、
元の文字列の中央から末尾（ゼロバイトでマーク）まで文字を印刷します。